\appendix

\chapter{Numerical simulation with axionbloch}

Key questions to be answered:
    \begin{itemize}
        \item 
        \item 
        \item 
    \end{itemize}
\begin{jointwork}
    ???
\end{jointwork}

\section{Hypothesis test}

Key questions to be answered:
    \begin{itemize}
        \item How to set up hypothesis? 
        \item 
        \item 
    \end{itemize}

\section{The Neyman--Pearson Lemma}

The \emph{Neyman--Pearson lemma} is a foundational result in statistical hypothesis testing. It establishes the form of the most powerful test for discriminating between two simple hypotheses at a fixed significance level.

\subsection{Problem Setup}

Suppose we observe data $x$ drawn from a probability distribution. We wish to test

\[
H_0 : x \sim f_0(x)
\]

against

\[
H_1 : x \sim f_1(x),
\]

where both $f_0(x)$ and $f_1(x)$ are fully specified probability density functions. Such hypotheses are referred to as \emph{simple hypotheses}.

The goal is to construct a statistical test that

\begin{itemize}
    \item maintains a fixed probability of false alarm (Type I error),
    \item while maximizing the probability of correctly detecting the alternative hypothesis (power).
\end{itemize}

\subsection{Likelihood Ratio}

The key quantity is the likelihood ratio

\[
\Lambda(x)
=
\frac{f_1(x)}{f_0(x)}.
\]

This ratio quantifies how much more likely the observed data is under the alternative hypothesis than under the null hypothesis.

Large values of $\Lambda(x)$ indicate that the observation is more consistent with $H_1$ than with $H_0$.

\subsection{Statement of the Neyman--Pearson Lemma}

For a fixed significance level $\alpha$, the most powerful test of $H_0$ against $H_1$ is given by

\[
\Lambda(x) > k,
\]

where the threshold $k$ is chosen such that

\[
P(\Lambda(x) > k \mid H_0) = \alpha.
\]

Equivalently, the optimal decision rule is:

\[
\text{Reject } H_0
\quad \text{if} \quad
\frac{f_1(x)}{f_0(x)} > k.
\]

The lemma therefore establishes that the likelihood-ratio test is optimal among all tests with the same false-positive rate.

\subsection{Type I Error and Power}

The significance level is defined as

\[
\alpha
=
P(\text{reject } H_0 \mid H_0 \text{ true}),
\]

which represents the probability of falsely rejecting the null hypothesis.

The power of the test is

\[
1-\beta
=
P(\text{reject } H_0 \mid H_1 \text{ true}),
\]

which measures the probability of correctly detecting the alternative hypothesis.

The Neyman--Pearson lemma guarantees that the likelihood-ratio test maximizes the power for a fixed $\alpha$.

\subsection{Gaussian Example}

Consider the case

\[
X \sim \mathcal{N}(\mu,\sigma^2),
\]

with known variance $\sigma^2$. We test

\[
H_0 : \mu = 0
\]

against

\[
H_1 : \mu = 1.
\]

Under the null hypothesis,

\[
f_0(x)
=
\frac{1}{\sqrt{2\pi\sigma^2}}
\exp\left(
-\frac{x^2}{2\sigma^2}
\right)\,,
\]

while under the alternative hypothesis,

\[
f_1(x)
=
\frac{1}{\sqrt{2\pi\sigma^2}}
\exp\left(
-\frac{(x-1)^2}{2\sigma^2}
\right).
\]

The likelihood ratio is therefore

\[
\Lambda(x)
=
\frac{f_1(x)}{f_0(x)}
=
\exp\left[
-\frac{(x-1)^2}{2\sigma^2}
+
\frac{x^2}{2\sigma^2}
\right].
\]

Expanding the exponent,

\[
(x-1)^2 = x^2 - 2x + 1,
\]

so that

\[
\log \Lambda(x)
=
\frac{x}{\sigma^2}
-
\frac{1}{2\sigma^2}.
\]

Hence, rejecting for large $\Lambda(x)$ is equivalent to rejecting for sufficiently large $x$.